\documentclass[11pt]{article}

\usepackage[a4paper,margin=1in]{geometry}
\usepackage{amsmath,amssymb,amsthm}
\usepackage{graphicx}
\usepackage[hidelinks]{hyperref}

\newtheorem{theorem}{Theorem}
\newtheorem{lemma}{Lemma}
\newtheorem{proposition}{Proposition}
\newtheorem{remark}{Remark}

\newcommand{\Hcal}{\mathcal H}
\newcommand{\N}{\mathbb N}
\newcommand{\ip}[2]{\left\langle #1,#2\right\rangle}
\newcommand{\norm}[1]{\left\lVert #1\right\rVert}

\title{A Block-Partition Criterion For Weighted Lower Semi-Frames}
\author{Automatic math-discovery smoke test}
\date{June 8, 2026}

\begin{document}
\maketitle

\section*{Source Problem}

The source paper is E. Zikkos,
\emph{The Gaussian Gabor system at the critical density is a weighted
lower semi frame}, arXiv:2606.00764, Functional Analysis.

It quotes Conjecture 3.1 of Balazs--Corso--Stoeva,
\emph{Weighted frames, weighted lower semi frames and unconditionally
convergent multipliers}, arXiv:2310.18957:

\begin{center}
\includegraphics[width=\linewidth]{figures/open_problem_crop.png}
\end{center}

Thus the ambient problem is whether every complete family in a separable
Hilbert space is a weighted lower semi-frame. The same passage notes that
the complete-minimal case is known and that the infinite-excess case
remains open; it also highlights the Muntz family
\(\{t^{\lambda_n}\}\subset L^2(0,1)\), with
\(\sum_n \lambda_n^{-1}=\infty\), as a natural test case.

\section*{Result Type}

This packet is a \textbf{partial result}. It gives a full affirmative
answer for the Muntz-family example in the quoted passage, and it proves
a general block-partition criterion, but it does not prove the full
complete-family conjecture.

\section*{Definitions And Convention}

A sequence \((g_n)\) in \(\Hcal\) is a lower semi-frame with lower bound
\(A>0\) if
\[
A\norm{f}^2\le \sum_n |\ip{f}{g_n}|^2
\qquad (f\in \Hcal).
\]
The sum is an extended nonnegative series. In particular, if the sum is
\(+\infty\) for every nonzero \(f\), then the lower semi-frame inequality
holds for every finite \(A>0\). A sequence \((f_n)\) is a weighted lower
semi-frame if \((\omega_n f_n)\) is a lower semi-frame for some positive
weights \(\omega_n\).

\section*{Block-Partition Criterion}

\begin{theorem}
Let \((f_n)_{n\in\N}\) be a sequence in a separable Hilbert space
\(\Hcal\). Suppose that
\[
\N=\bigsqcup_{j=1}^{\infty} I_j
\]
and that each subfamily \((f_n)_{n\in I_j}\) is complete in \(\Hcal\).
Then \((f_n)\) is a weighted lower semi-frame. More precisely, for every
\(A>0\) there are positive weights \((\omega_n)\) such that
\[
\sum_n |\ip{f}{\omega_n f_n}|^2=+\infty
\]
for every nonzero \(f\in\Hcal\).
\end{theorem}

\begin{proof}
Fix an orthonormal basis \((e_m)\) of \(\Hcal\). Since the \(j\)-th block
is complete, choose a finite subset \(G_j\subset I_j\) such that
\[
\operatorname{dist}(e_m,\operatorname{span}\{f_n:n\in G_j\})<\frac1j
\qquad (1\le m\le j).
\]
Let
\[
E_j=\operatorname{span}\{f_n:n\in G_j\}
\]
and let \(P_j\) be the orthogonal projection onto \(E_j\). The preceding
approximation property implies \(P_j\to I\) strongly: first check this
on finite linear combinations of basis vectors, then use uniform
boundedness of the projections and density.

For each \(j\), the finite family \((f_n)_{n\in G_j}\) spans \(E_j\).
Therefore
\[
\delta_j
=\inf_{\substack{x\in E_j\\ \norm{x}=1}}
\sum_{n\in G_j}|\ip{x}{f_n}|^2
\]
is strictly positive. Define
\[
\omega_n=\left(\frac{A}{\delta_j}\right)^{1/2}
\qquad (n\in G_j),
\]
and assign any positive value, say \(1\), to indices not belonging to
\(\bigcup_jG_j\). The sets \(G_j\) are pairwise disjoint because the
blocks \(I_j\) are disjoint.

For every \(f\in\Hcal\),
\[
\sum_{n\in G_j}|\ip{f}{\omega_n f_n}|^2
=\frac{A}{\delta_j}\sum_{n\in G_j}|\ip{P_jf}{f_n}|^2
\ge A\norm{P_jf}^2 .
\]
If \(f\ne0\), then \(\norm{P_jf}\to\norm{f}>0\). Hence the block
contributions are bounded below by a positive constant for all large
\(j\), and summing over the disjoint blocks gives
\[
\sum_n |\ip{f}{\omega_n f_n}|^2=+\infty .
\]
\end{proof}

\section*{Application To The Muntz Family}

\begin{lemma}
Let \((a_n)\) be a positive sequence with \(\sum_n a_n=\infty\). Then
\(\N\) can be partitioned into countably many disjoint sets \(J_k\) such
that \(\sum_{n\in J_k}a_n=\infty\) for every \(k\).
\end{lemma}

\begin{proof}
Split \(\N\) into consecutive finite intervals \(B_m\) such that
\(\sum_{n\in B_m}a_n\ge1\) for every \(m\). This is possible because
every tail of the positive divergent series diverges. Choose a map
\(s:\N\to\N\) such that \(s^{-1}(k)\) is infinite for every \(k\), and
set
\[
J_k=\bigcup_{m:s(m)=k}B_m .
\]
These sets are disjoint, cover \(\N\), and each corresponding subseries
dominates an infinite sum of ones.
\end{proof}

\begin{proposition}
Assume \(\lambda_n\in\N\) and \(\sum_n\lambda_n^{-1}=\infty\). Under the
Muntz--Szasz completeness criterion used in arXiv:2606.00764, the family
\[
\{t^{\lambda_n}\}\subset L^2(0,1)
\]
is a weighted lower semi-frame for \(L^2(0,1)\).
\end{proposition}

\begin{proof}
Apply the lemma to \(a_n=\lambda_n^{-1}\). The resulting partition
\(\N=\bigsqcup_kJ_k\) satisfies
\[
\sum_{n\in J_k}\lambda_n^{-1}=\infty
\qquad (k\ge1).
\]
By the Muntz--Szasz criterion as invoked in the source paper, each
subfamily \(\{t^{\lambda_n}:n\in J_k\}\) is complete in \(L^2(0,1)\).
The block-partition theorem now gives positive weights making the full
Muntz family a lower semi-frame.
\end{proof}

\section*{Why This Is Not A Full Solution}

The argument requires a partition of the index set into countably many
complete subfamilies. That condition holds for the Muntz example because
the divergent series \(\sum_n\lambda_n^{-1}\) can be split into
countably many divergent subseries. It is not shown here, nor in the
source paper, that every complete family in the remaining infinite-excess
regime admits such a partition.

Thus the result should remain in the \texttt{partial} directory. It would
become a full solution of the complete-family conjecture only after
removing the block-partition hypothesis or proving that hypothesis for
all complete families covered by the conjecture.

\section*{Human Review Recommendation}

The proof is short and appears internally consistent. The main points to
check are the lower semi-frame convention in the two cited papers, the
finite-dimensional lower bound \(\delta_j>0\), and the exact
Muntz--Szasz completeness statement used for subfamilies of
\(\{t^{\lambda_n}\}\).

\end{document}
